\subsection*{Chapter 1: The Landscape}

\paragraph{Exercise 1.1}
By the definition of conditional probability, $p(x_1,\dots,x_T)=p(x_T\mid x_{<T})\,p(x_1,\dots,x_{T-1})$; applying
this recursively down to $p(x_1)$ gives the product. No modelling assumption is involved, so an exact
autoregressive model can represent any distribution. With a causal mask, the output at position $t$ depends
only on $x_{\le t}$, so position $t$ is a valid prediction of $x_{t+1}$ for every $t$: one forward pass over
$T$ tokens yields $T$ supervised predictions (the last uses the next token in the stream). A masked LM
computes a loss only at the $\approx0.15T$ masked positions, because unmasked positions can see their own
answer through bidirectional attention.

\paragraph{Exercise 1.2}
Perplexity is $e^{2.0}\approx7.39$: the model is, on average, as uncertain as if it were choosing uniformly
among about 7.4 tokens. A uniform model has loss $\ln 50{,}257\approx10.82$ nats and perplexity 50{,}257.
This is a useful sanity check: a freshly initialised model should start near 10.8.

\paragraph{Exercise 1.3}
Wikipedia: $0.03\times300\text{B}=9\text{B}$ tokens sampled from a 3B-token source, about 3 epochs.
Common Crawl: $0.60\times300\text{B}=180\text{B}$ from 410B, about 0.44 epochs. Repeating small, clean,
information-dense sources raises the quality of the effective training distribution. The cost is a risk of
memorisation and diminishing returns after a few epochs, so heavy repetition is reserved for the
highest-quality data.

\paragraph{Exercise 1.4}
$C=64\times30\times86{,}400\,\text{s}\times989\times10^{12}\times0.35\approx5.74\times10^{22}$ FLOPs. With
$C=6ND$ and $D=20N$, $C=120N^2$, so $N=\sqrt{C/120}\approx2.2\times10^{10}$ (about 22B parameters) and
$D\approx4.4\times10^{11}$ tokens. If the model will be served heavily, total cost is training plus
$2N$ FLOPs per served token. Shrinking $N$ (say to 7B) and spending the same compute on more tokens
($D=C/(6N)\approx1.4$T) gives a slightly worse model at equal compute but a $3\times$ cheaper one to serve, which
is usually the better trade.

\paragraph{Exercise 1.5}
bf16: weights $4.02\times10^9\times2=8.04$\,GB, KV at 4{,}096 tokens $=147{,}456\times4{,}096=0.60$\,GB, total
8.64\,GB, giving $8.64/3.35\times10^{12}\approx2.58$\,ms, or about 387 tokens/s. int4 weights: 2.01\,GB + 0.60\,GB =
2.61\,GB, about 0.78\,ms, or about 1{,}280 tokens/s. Weights still dominate, but the KV share of traffic rises from
7\% to 23\%; at longer contexts or larger batches the KV cache becomes the dominant term, which motivates
KV-cache quantisation (\cref{ch:efficient}).

\paragraph{Exercise 1.6}
Per step, compute time is $B\cdot2N/F$ and memory time is $Nb/W$ (weights read once, shared by the batch).
Equating gives $B^\star=bF/(2W)$. For an H100 in bf16, $B^\star=2\times989/(2\times3.35)\approx295$, matching the
lab's 296--320 (the lab includes the LM head and attention FLOPs and embedding bytes, which shift the
value slightly). With long contexts each sequence adds its own KV bytes, which are \emph{not} shared across
the batch. Attention over the cache has intensity of about 1\,FLOP/byte regardless of batch size, so the step
needs a larger batch, or never becomes compute-bound when KV traffic dominates.

\paragraph{Exercise 1.7}
Qwen3-8B uses 147{,}456 bytes of KV per token. Short answers: $500+300=800$ tokens, 118\,MB per request, about
339 concurrent requests in 40\,GB. With thinking: about $500+6{,}000=6{,}500$ tokens, 958\,MB per request, about 41
concurrent requests. Under memory-bound decoding, throughput scales roughly with batch, so throughput in
tokens/s drops by about $8\times$. Each request also emits about $8\times$ more tokens, so requests per second fall by
roughly $60\times$.

\paragraph{Exercise 1.8}
$0.97^{30}\approx0.40$. If wrong answers are spread over, say, ten distinct values, a simulation of plurality
voting (counting ties as failures) gives about 0.57 success at 5 samples, 0.81 at 11, 0.89 at 15 and 0.96 at 21,
so 15--17 samples are needed for 0.9. Independence is optimistic because a model's errors are correlated: the
same misconception recurs across samples, concentrating wrong votes on one value and making plurality voting
far less effective.

\paragraph{Exercise 1.9}
Example probes: (i) ``All squares have five sides. ABCD is a square. How many sides does ABCD have?''
(closed-world answer 5; a good reasoner flags the false premise). (ii) A freshly generated puzzle with random
names and quantities, for instance a scheduling constraint problem produced by a script with a computed
answer. (iii) A question whose answer runs against frequency, for instance ``Which is heavier, a kilogram of
feathers or 900\,g of steel?''. Base models often continue the text without answering; instruction-tuned models
answer fluently but fall for (iii) more often; reasoning models do best on (ii). Score automatically with
programmatic answers and the verifier of \cref{ch:eval}, and for (i) with a rubric for ``flags the
inconsistency'' judged by an LLM (\cref{ch:eval}).

\paragraph{Exercise 1.10}
A reasonable ranking is distillation, then RL, then inference-time scaling. Strict latency rules out
heavy per-query sampling. One week of compute is enough to fine-tune on traces from a stronger teacher,
which is cheap and stable, and the fine-tuned model serves at its normal speed. RL is attractive if a
reliable automatic reward exists (for instance, review comments that can be checked against later bug fixes)
but is costlier to make work; it can follow distillation as a second stage. A light form of inference-time
scaling (one short chain of thought) can still be used if the latency budget allows.
